\documentclass[12pt, ukrainian]{article}
\usepackage[a4paper]{geometry}
\setlength\parindent{0mm}
\geometry{top=1.0cm,bottom=1.0cm,left=1.3cm,right=1.3cm}
\pagestyle{empty}
\usepackage[T2A]{fontenc}
\usepackage[utf8]{inputenc}
\usepackage[ukrainian]{babel}
\usepackage{graphicx}
\usepackage{caption}
\usepackage{caption}
\DeclareCaptionFormat{nocaption}{}
\captionsetup{format=nocaption,aboveskip=0pt,belowskip=0pt}
\usepackage[Export]{adjustbox}
\adjustboxset{max size={0.9\linewidth}{0.9\paperheight}}
\usepackage{verbatim, listings, clrscode3e, fancyvrb}
\def\code#1{\Verb!#1!}
\begin{document}
%begin
\begin {large}

\textbf{\_\_.11.2025 Контрольна робота, варіант  \No { 9 }}\par


{

\begin{center}
\textbf{Завдання 1}
\end{center}
Нехай int var=13; Один потік збільшив змінну var на 3, а інший збільшив її в три рази.

гонка даних (data race), стан гонки (race condition)

Виберіть усі правильні твердження (якщо такі є).\par
\vspace{2mm}
(\,1\,) Тут гарантовано є стан гонки.%good

(\,2\,) Тут гарантовано є і стан гонки, і гонка даних.%bad

(\,3\,) Тут може не бути стану гонки.%bad

(\,4\,) Тут може бути тільки гонка даних.%bad



\begin{center}
\textbf{Завдання 2}
\end{center}
Виберіть усі правильні твердження (якщо такі є).\par
\vspace{2mm}

(\,1\,) Функція async запускає виконання в новому потоці.%bad

(\,2\,) Тільки від компілятора залежить, чи запустить функція async виконання в новому потоці.%bad

(\,3\,) Довільний потік може розблокувати заблокований м’ютекс.%bad

(\,4\,) Порядок обчислень у потоці для одних тих самих вхідних даних може залежати від використовуваного компілятора.%good



\begin{center}
\textbf{Завдання 3}
\end{center}
Виберіть усі правильні твердження (якщо такі є).\par
\vspace{2mm}

(\,1\,) Вільна від блокувань структура даних гарантує, що жоден потік не буде очікувати доступу до виконання операцій над цією структурою даних.%bad

(\,2\,) Операція над вільною від блокувань структурою даних завжди виконається за обмежений час.%bad

(\,3\,) Кожна вільна від очікувань структура даних є структурою даних, вільною від блокувань.%good

(\,4\,) Для уникнення гонки даних структура даних може використовувати не більше одного м’ютекса.%bad



\begin{center}
\textbf{Завдання 4}
\end{center}
Виберіть усі правильні твердження (якщо такі є).\par
\vspace{2mm}

(\,1\,) Потік може мати доступ до автоматичних змінних іншого потоку, але це небезпечно.%good

(\,2\,) Потік має доступ тільки до своїх власних даних та до глобальних статичних даних.%bad

(\,3\,) Потік мови С++ не може отримати доступ до автоматичних змінних іншого потоку.%bad

(\,4\,) Автоматичні змінні не можуть розділятися між потоками.%bad



\begin{center}
\textbf{Завдання 5}
\end{center}
У цьому фрагменті коду
\begin{verbatim}
template <class T> void f(T a){}

void g(){
    f(unique_ptr<int>(new int(3)))+(unique_ptr<int>(new int(35)));
\end{verbatim}

(\,1\,) є гонка даних%bad

(\,2\,) є стан гонки%bad

(\,3\,) є невизначена поведінка%bad

(\,4\,) можливий виток ресурсів%good


}
\end {large}


\end{document}

